%% Teste completo de integração bibliográfica
%% Este arquivo demonstra todos os recursos de bibliografia da classe UnB-PPGT

\documentclass[mestrado]{UnB-PPGT}

%% Metadados mínimos para teste
\titulo{Demonstração Completa do Sistema Bibliográfico}
\autor{Teste}{Usuario}
\orientador{Prof. Dr. Orientador Teste}{Universidade de Brasília}
\diamesano{1}{1}{2024}
\membrobancafunc{Prof. Dr. Orientador Teste}{Universidade de Brasília}{Orientador}
\membrobancafunc{Prof. Dr. Membro 1}{Universidade de Brasília}{Examinador}
\membrobancafunc{Prof. Dr. Membro 2}{Universidade Federal de Goiás}{Examinador Externo}

\begin{document}

\pretextual

\begin{resumo}
Este documento demonstra o sistema bibliográfico integrado da classe UnB-PPGT, incluindo diferentes tipos de citação, formatação automática de referências e compatibilidade com natbib.
\end{resumo}

\textual

\chapter{Demonstração do Sistema Bibliográfico}

Este capítulo demonstra os recursos bibliográficos implementados na classe UnB-PPGT.

\section{Citações Básicas}

\subsection{Citações no Texto}
Exemplo de citação integrada ao texto: \citeonline{vasconcellos2012} apresenta uma análise abrangente da mobilidade urbana no contexto brasileiro.

\subsection{Citações entre Parênteses}
A mobilidade urbana é um tema complexo que envolve múltiplas dimensões \cite{vasconcellos2012}.

\section{Citações Avançadas}

\subsection{Citações com Página}
Para citações específicas: \citepag{p. 45}{silva2019} discute os padrões de mobilidade.

Citação no texto com página: \citeonlinepag{p. 123}{oliveira2021} aborda questões de sustentabilidade.

\subsection{Citações Múltiplas}
Vários autores tratam do tema de mobilidade urbana \cite{vasconcellos2012,silva2019,oliveira2021}.

\subsection{Citações de Autor e Ano}
Segundo \citeauthor{costa2018}, em \citeyear{costa2018}, o planejamento urbano é fundamental.

Formato autor-ano: \citeauthoryear{ferreira2021}.

\section{Citações Especializadas}

\subsection{Citação Apud}
Exemplo de citação indireta: \citeapud{costa2018}{silva2019}.

\subsection{Citações de Legislação}
Referência à legislação: \citelaw{brasil2012}.

\subsection{Citações de Dados e Relatórios}
Dados estatísticos: \citedata{ibge2022}.

Relatórios técnicos: \citereport{detran2023}.

\section{Tipos de Referências}

\subsection{Livros}
Referência a livro: \cite{yin2015}.

\subsection{Artigos em Periódicos}
Artigo científico: \cite{silva2019}.

\subsection{Teses e Dissertações}
Trabalho acadêmico: \cite{santos2020}.

\subsection{Recursos Online}
Fonte online: \cite{who2021}.

\subsection{Normas Técnicas}
Norma ABNT: \cite{abnt2020}.

\chapter{Formatação da Bibliografia}

A bibliografia é formatada automaticamente seguindo as normas acadêmicas, com:

\begin{itemize}
\item Espaçamento adequado entre entradas
\item Fonte menor para melhor legibilidade
\item Formatação consistente para diferentes tipos de referência
\item Integração automática ao sumário
\item Suporte a hyperlinks (quando habilitado)
\end{itemize}

\postextual

%% Bibliografia - demonstra o comando aprimorado
\bibliography{bibliografia}

\end{document}